\section{Local operators and amplification}\label{sec:local-amplification}

The reflected perturbation contains a number of terms proportional to
the volume. To control one local insertion uniformly, we repeat it
in large diagonal cells and then fold those cells in coordinate
directions. The coordinate trace bounds each repeated operator by
its local norm. Taking the root corresponding to the number of
copies leaves an error that tends to one when volume tends to
infinity before the perturbation is removed. The local norm, insertion
bound and amplification are the constructions of
\cite[Lemma~8.13, Proposition~8.26 and Theorem~8.32]{ChamberQED}.

The estimate to be proved is
\[
 |\mu_i^t(G)|\le C\gamma_{t,i}(k),\qquad
 \gamma_{t,i}(k)=
 \exp\bigl(2t|O|_1k^3+|Q|k^3\Delta(m_i)\bigr).
\]
Here $C$ is the coordinate bound of the local insertion $G$, defined
in \eqref{eq:diag-coordinate-bound}; for the insertion $G_X$ of
\eqref{eq:diag-operator-insertions} it is $\|X\|$.
The integer $k$ depends on the fixed support
radius, cell width and side sequence, but not on $t$, $i$ or the
location of the support. We choose it before taking limits. The
pressure error $\Delta(m_i)$ tends to zero, so
$\lim_{t\downarrow0}\lim_{i\to\infty}\gamma_{t,i}(k)=1$.
\Cref{thm:diag-amplification} gives the precise assertion and proof.
We first construct the local norm and the insertions to which it applies.

\subsection{Local operators at a coordinate link plane}\label{subsec:diag-local-operators}

An insertion must carry an operator norm that is uniform in the
cross-section and in electric characters. Fix a finite set $S$ of
links of $\Z^2$ and a nonempty finite set $\Lambda$ of sites
containing their endpoints. Let $\mathcal F_\Lambda$ be the Fock
space of crossing modes over $\Lambda$, and put
\[
 H_S=L^2(\R^S)\otimes\mathcal F_\Lambda,
 \qquad E_b=-i\partial_{A(b)}.
\]
The diagonal unitary $u(a)$ on $\mathcal F_\Lambda$ has integer mode
charges. An operator $\Upsilon$ has charge $z\in\Z^\Lambda$ when
$u(a)\Upsilon u(a)^{-1}=e^{iq\langle a,z\rangle}\Upsilon$.

Let $\mathfrak A(S)$ be the span of the functions
$e^{-Q(\xi_{S'})+\langle\lambda,\xi_{S'}\rangle}$, where
$S'\subset S$, $Q$ is real and positive definite on $\R^{S'}$,
and $\lambda\in\C^{S'}$. The empty subset gives the constant one.
These functions are bounded and uniformly continuous. Their span
is closed under products, conjugation and real translations: products
add the positive quadratic forms on the union of their supports,
while translation completes the same quadratic exponent.

For $\delta\in\R^S$, extended by zero to the full slice, require
$d_s^*\delta=qz(\delta)$ with $z(\delta)\in\Z^\Lambda$. The local
operator algebra consists of the finite sums
\begin{equation}\label{eq:diag-operator-normal-form}
 X=\sum_\delta e^{i\langle\delta,A\rangle}\Upsilon_\delta(E_S),
 \qquad\Upsilon_\delta\in\mathfrak A(S)\otimes
                      \operatorname{End}_{z(\delta)}(\mathcal F_\Lambda).
\end{equation}
We denote it by $\mathfrak W(S,\Lambda)$. The Weyl identity gives
\begin{align}\label{eq:diag-operator-rules}
 (e^{i\delta A}\Upsilon(E_S))(e^{i\delta'A}\Upsilon'(E_S))
       &=e^{i(\delta+\delta')A}\Upsilon(E_S+\delta')\Upsilon'(E_S),
          \notag\\
 (e^{i\delta A}\Upsilon(E_S))^*
       &=e^{-i\delta A}\Upsilon^*(E_S-\delta).
\end{align}
Here $\delta A$ abbreviates the real pairing. These formulas make
$\mathfrak W(S,\Lambda)$ a unital $*$-algebra. The normal form is
unique. In Fourier variables its terms send a function supported
near $\xi_0$ to functions supported near $\xi_0+\delta$; small
enough neighborhoods separate the finitely many distinct exponents.
If the sum vanishes as an operator, every continuous symbol vanishes
on each such neighborhood and hence everywhere.

Every $X$ commutes with the local gauge transformations
$\Phi(A)\mapsto u(a)\Phi(A-d_sa)$. The phase change
$e^{-i\langle\delta,d_sa\rangle}$ cancels the Fock charge phase.
The total charge of $z(\delta)$ is zero, since it is the divergence
of a finitely supported cochain. Each Fock term is therefore even.
For $\beta\in\R^S$, define
\begin{equation}\label{eq:diag-alpha}
 \alpha_\beta X=e^{i\langle\beta,A\rangle}Xe^{-i\langle\beta,A\rangle}
    =\sum_\delta e^{i\langle\delta,A\rangle}
                                      \Upsilon_\delta(E_S-\beta).
\end{equation}
This is an isometric $*$-automorphism on $H_S$.

For a rectangular cross-section $\Xi$ containing the smallest integer
rectangle spanned by $\Lambda$,
the same formulas define $X_\Xi$ on the real slice quotient space
$H^l$ of \cref{subsec:slab-coordinate-trace}. The Fock operators
are localized into the full crossing Fock space by $\iota$.
The localization is the tensor inclusion for the exterior space, with the reordering
sign and cross-section character attached to each representative;
it is a unital injective $*$-homomorphism. The phase
$e^{i\langle\delta,A\rangle}$ is well defined on the quotient.
Indeed, on the support hull an integer closed one-cochain is an
integer coboundary, by the box Poincare argument. Thus for
$\gamma\in\Gamma_s$,
$\langle\delta,\gamma\rangle\in2\pi\Z$, using
$d_s^*\delta=qz(\delta)$. The local Weyl formulas therefore remain
valid on the quotient. Uniqueness of the normal form makes the
representation independent of its presentation. It commutes with
$P_G$ and with functions of the electric variables outside $S$.

Define the norm needed for characters by
\begin{equation}\label{eq:diag-sharp-norm}
 \|X\|_\#=\sup_{\Xi,\beta}\| (\alpha_\beta X)_\Xi\|.
\end{equation}
The support hull must embed in the cross-section; the supremum is
over such cross-sections and all $\beta\in\R^S$. The electric spectrum
of a quotient can be smaller than $\R^S$, so its unshifted operator
norm need not equal the local norm.

\begin{proposition}\label{prop:diag-local-norm}
For $X\in\mathfrak W(S,\Lambda)$,
\[
 \|X\|_\#=\|X\|_{\mathcal L(H_S)}.
\]
The norm is unchanged by enlarging $S,\Lambda$. In particular it is
a $C^*$-norm, $\|1\|_\#=1$, and every quotient realization satisfies
$\|(\alpha_\beta X)_\Xi\|\le\|X\|_\#$.
\end{proposition}

\begin{proof}
First remove the Fock localization. Its tensor extension is an
injective $*$-homomorphism and hence isometric; we may work on
$L^2(\R^{\Sigma^{(1)}}/\Gamma_s)\otimes\mathcal F_\Lambda$ with the
original finite Fock operators. Call this realization $X_\Xi^0$.
Introduce an additional variable $z\in\R^S$, with electric operators
$E'=-i\partial_z$, on
\[
 L^2(\R^S_z)\otimes
 L^2(\R^{\Sigma^{(1)}}/\Gamma_s)\otimes\mathcal F_\Lambda.
\]
The operator $X(z,E')$ has norm $\|X\|$, since it is $X$ tensored
with the identity. Choose a Borel representative $A^\flat$ of the
quotient and use the two unitaries
\[
 (U_1\Phi)(z,[A])=\Phi(z+A^\flat|_S,[A]),\qquad
 (U_2\Phi)(z,[A])=\Phi(z,[A-z]),
\]
where $z$ is extended by zero outside $S$. The first conjugation
multiplies $e^{i\delta z}$ by $e^{i\delta A}$. The second replaces
$E'$ by $E'+E_S$ and cancels the factor $e^{i\delta z}$. Hence
\[
 U_2U_1X(z,E')(U_2U_1)^{-1}
       =\sum_\delta e^{i\delta A}\Upsilon_\delta(E_S+E').
\]
Fourier transform in $z$. This becomes a direct integral of
$T_p=(\alpha_{-p}X)_\Xi^0$, $p\in\R^S$. Uniform continuity of the
symbols makes $p\mapsto T_p$ norm continuous. Its direct-integral
norm is therefore $\sup_p\|T_p\|$: the reverse inequality follows
by using a unit vector nearly maximizing $T_{p_0}$ on a small ball
about $p_0$. It follows that
\[
 \|X\|=\sup_p\| (\alpha_{-p}X)_\Xi^0\|
        =\sup_p\| (\alpha_{-p}X)_\Xi\|.
\]
This holds for each admissible cross-section and proves the asserted
supremum identity. Enlarging $S$ tensors with an identity on the
additional $L^2$ factor; enlarging $\Lambda$ uses the same isometric
Fock inclusion. Both preserve the norm.
\end{proof}

\subsection{Insertions and a coordinate trace bound}\label{subsec:diag-insertions}

Let $c=j+\tfrac12$. Write $\operatorname{Ins}_c$ for the crossing
insertion map, and $\omega$ for its diagonal mode weights. For an
even matrix kernel $k(A',A)$ on $\R^S$, set
\begin{equation}\label{eq:diag-kernel-insertion}
 \begin{split}
 G[k](\widehat A)={}&\operatorname{Ins}_c\!\left[
   \omega^{1/2}k(\widehat A_{j+1}|_S,
           (\widehat A_j+d_s\widehat a_j)|_S)
             u(\widehat a_j)\omega^{1/2}\right]\\
 &\quad\cdot\prod_{b\in S}g_1(F_j^{\mathrm{el}}(b))^{-1}
                      \exp(S_c^\times(\Lambda)).
 \end{split}
\end{equation}
Here $S_c^\times(\Lambda)$ is the part of the fermion action crossing
the plane over $\Lambda$. The two factors on the second line remove
the original electric and fermion crossing weights on the support
of $k$. Thus $G[k]$ inserts $k$ in their place. Its definition uses
the infinite lattice and is independent of a torus.

We require
\begin{align}\label{eq:diag-kernel-conditions}
 u(a)^{-1}k(A'+d_sa,A+d_sa)u(a)&=k(A',A),\notag\\
 \|k(A',A)\|&\le C\exp(-\epsilon|A'-A|^2),\qquad\epsilon>0.
\end{align}
Choose a box containing the two copies of $S$, the temporal links
over $\Lambda$, and the unit cubes centered at $(j,y)$ and $(j+1,y)$,
$y\in\Lambda$. In what follows we use the cube $B^\#(c,S,\Lambda)$
obtained by taking twice the largest half-width of the smallest
rectangular box containing these sets, with the same center. This
choice is covariant under coordinate translations and reflections.
The insertion reads only these links and the crossing fields, so
it is supported in $B^\#$.

The insertion is an even, gauge invariant local integrand. Indeed,
an even kernel has entries only between occupation vectors of the
same parity, so every monomial
$\Theta_c\zeta_J\,\zeta_I$ in its insertion is even. If the gauge
parameter is $\lambda$, its upper and lower slice values change
the matrix in \eqref{eq:diag-kernel-insertion} to
\[
 u(\lambda_{j+1})^{-1}
   k(\widehat A_{j+1},\widehat A_j+d_s\widehat a_j)
       u(\widehat a_j)u(\lambda_j).
\]
The matrix entry acquires precisely the inverse of the phase of
$\Theta_c\zeta_J\,\zeta_I$. The electric curvature and the crossing
action are gauge invariant. Finally, the kernel estimate and
$g_1(x)^{-1}=\sqrt{2\pi}e^{x^2/2}$ bound every Grassmann coefficient
by
\begin{equation}\label{eq:diag-insertion-growth}
 C'\exp\left((\tfrac12-\epsilon')
                 \sum_{b\in S}F_j^{\mathrm{el}}(b)^2\right),
 \qquad\epsilon'=\min(\epsilon,\tfrac12).
\end{equation}
The other factors have bounded coefficients. This is the local
integrand bound of \cref{subsec:slab-local}, since its electric
plaquettes lie in the box.

For $X\in\mathfrak W(S,\Lambda)$ let $k_X^{(s)}$ be the continuous
kernel of $e^{-s|E_S|^2}Xe^{-s|E_S|^2}$, with $s=\tfrac14$ or
$\tfrac18$, and write
\begin{equation}\label{eq:diag-operator-insertions}
 G_X=G[k_X^{(1/4)}],\qquad G'_X=G[k_X^{(1/8)}].
\end{equation}
Both kernels satisfy \eqref{eq:diag-kernel-conditions}. To see their
decay, consider one normal-form term $e^{i\delta A}\Upsilon(E)$ with
$\Upsilon(\xi)=e^{-Q(\xi)+\lambda\xi}\Upsilon_0$. Identify $Q$
with its symmetric matrix, so $Q(\xi)=\langle\xi,Q\xi\rangle$,
extending $Q$ by
zero to the other links. With $Q_s=4sI+2Q$ and
$\lambda_s=\lambda-2s\delta$, its kernel is
\begin{equation}\label{eq:diag-regularized-kernel}
 \begin{split}
 k_X^{(s)}(A',A)={}&e^{i\delta A'}(2\pi)^{-|S|/2}
          (\det Q_s)^{-1/2}e^{-s|\delta|^2}\\
 &\cdot\exp\left(\tfrac12
  \langle\lambda_s+i(A'-A),Q_s^{-1}
                  (\lambda_s+i(A'-A))\rangle\right)\Upsilon_0.
 \end{split}
\end{equation}
The real part of this exponent has a strictly negative quadratic
form in $A'-A$, with a fixed linear term. It gives the required
Gaussian decay after decreasing its exponent. Gauge covariance
follows by conjugating the regularized operator by the gauge
unitary; equality of their continuous kernels gives the first
identity in \eqref{eq:diag-kernel-conditions}. Finite sums preserve
these properties. The identity operator has kernel
$\prod_{b\in S}g_1(A'(b)-A(b))$ when $s=\tfrac14$, so
$G_1=1$; unused links and modes cancel in the same way.

Reflection acts on these insertions by adjoints:
\begin{equation}\label{eq:diag-insertion-reflection}
 \Theta_cG_X=G_{X^*},\qquad\Theta_cG'_X=G'_{X^*}.
\end{equation}
The kernel is changed to $k(A,A')^*$. At the same time the two
slices exchange, $\widehat a_j$ changes sign, and covariance moves
$u(\widehat a_j)$ to the required side of the kernel. This gives
\eqref{eq:diag-insertion-reflection}, including its field signs.
We also need coordinate translations and spatial reflections to
preserve the operator norm, since each coordinate fold contains
such images.

For a spatial reflection $\theta$ in the second or third coordinate,
the spin matrix exchanges the two eigenspaces of $\gamma_1$.
Its unitary matrix $v$ on crossing modes preserves their weights
and reverses their charges. On exterior products the reordering
sign is $r_I=(-1)^{|I|(|I|-1)/2}$; hence the Fock unitary is
$U=(\bigwedge v)r$. The anticommutation of the two field reflection
maps contributes the sign $(-1)^{\deg}$, and moving the two mode
products past one another contributes $(-1)^{|I||J|}$. For the
parity-matched entries of an even kernel these signs leave exactly
$r_Ir_J$. Consequently
\[
 \Theta_\theta\operatorname{Ins}_c[\Upsilon(\widehat A)]
 =\operatorname{Ins}_c[
       U\overline{\Upsilon(\theta^*\widehat A)}U^*].
\]
The antiunitary
$(J\Phi)(A)=U\overline{\Phi(\theta_s^*A)}$ gives this kernel
transformation. It sends an oriented electric operator to the
negative of the reflected electric operator, and sends
$e^{i\delta A}$ to $e^{-i\langle\theta_s^*\delta,A\rangle}$.
Thus it preserves the charge constraint, the local algebra and
the heat factors $e^{-s|E|^2}$. Translations act by the corresponding
unitary relabeling. Together with adjoints at the link plane, these
facts show that each translate or signed coordinate reflection
of $G_X$ is an insertion $G_{X^g}$ with
$\|X^g\|=\|X\|$. The same holds for $G'_X$. The box $B^\#$
transforms with the insertion.

An even local integrand $G$ in a box $B$ satisfies the coordinate
bound $C\ge0$ if, on every rectangular coordinate torus
$T=(\Gamma,\chi)=C(n_1,n_2,n_3)$ with $n_\mu\ge4$,
for every nonempty finite family of images $g_aB$ under lattice
translations and signed coordinate reflections that embeds
disjointly in $\Gamma$, and every $\beta\in C^2(\Gamma)$,
\begin{equation}\label{eq:diag-coordinate-bound}
 \left|Z_T\left[\varphi_\beta\prod_{a=1}^rG^{g_a}\right]\right|
       \le C^rZ_T[1],\qquad
       \varphi_\beta(F)=e^{i\langle\beta,F\rangle},
\end{equation}
provided that, in each coordinate direction $\mu$, there is a
plaquette $\sigma$ orthogonal to $e_\mu$ such that none of the
boxes or their period translates contains any plaquette
$\sigma+k e_\mu$, $k\in\Z$. This condition places the flux
representative used in the slice trace on plaquettes read by
none of the insertions.
When the coordinate cell structure cuts $T$ and its walls are at
half-integers, boxes contained in its open cells satisfy the
free-column condition. Choose a plaquette orthogonal to the chosen
direction that meets a wall normal to another direction.
Period translations permute the open cells, while every plaquette
in this column meets that wall. Thus none of these plaquettes
is contained in a box or any of its period translates.

\begin{proposition}\label{prop:diag-coordinate-insertions}
The insertion $G_X$ satisfies \eqref{eq:diag-coordinate-bound} with
$C=\|X\|$. The insertion $G'_X$ satisfies it with
$C=2^{|S|/2}\|X\|$. A bounded bosonic integrand satisfies it with
$C=\|G\|_\infty$. These bounds survive box enlargement,
multiplication by a curvature character, and transport by a
signed permutation of the coordinates.
\end{proposition}

\begin{proof}
We give the trace calculation, including the support condition.
Choose the time direction to be the first coordinate. By the
covariance just proved, the images of the insertions have the same
operator norms and are at link planes $j+\tfrac12$. The free column
lets us choose the magnetic flux plaquette $\sigma_0$ so that no
plaquette in its column is contained in an insertion box.
At a fixed step $j$ their spatial supports and modes
are disjoint. Let $S_j$ be the union of the supports of the primed
insertions at that step.

The unfolded path variables are genuine local lifts on each
$[j,j+1]\times\operatorname{hull}(\Lambda_a)$. On electric
plaquettes their coboundary is $A_{j+1}-A_j-d_sa_j$. On magnetic
plaquettes it is $d_sA_j$; the flux correction is supported at
$\sigma_0$, which is not a plaquette of the local hull.
Two local lifts with
these curvatures and phases differ by an integer exact cochain
$\frac{2\pi}{q}d\lambda$. Gauge invariance therefore allows the
insertion to be read directly at the path variables. This is why
the free-column hypothesis is needed.

At each crossing, the factors $e^{S_c^\times(\Lambda_a)}$ remove
the crossing weights on the local modes. The remaining crossing
factors, together with their insertions, form
\[
 \operatorname{Ins}_c\left[
   \omega^{1/2}\left(\prod_a\iota(k_a(A_{j+1},A_j+d_sa_j))\right)
                   u(a_j)\omega^{1/2}\right].
\]
The even localized matrices commute because their mode sets are
disjoint. The factors $g_1^{-1}$ similarly remove the electric
Gaussian densities on the local links. The slice determinant
factorization of \cref{subsec:slab-coordinate-trace} now gives
the path product with these modified crossing kernels. Its
Gaussian majorant is integrable: every modified kernel has
strict decay in its electric difference, all remaining links
retain $g_1$, and the magnetic multiplier has its Gaussian
weight. Hence unfolding, periodization and the path trace are
valid for this product.

Separate the curvature character into electric factors
$e^{i\langle\beta_j,F_j^{\mathrm{el}}\rangle}$ and magnetic
unitaries $u_j$. Multiplying a regularized kernel by
$e^{i\langle\beta_j,A'-A\rangle}$ changes the corresponding
operator to
\[
 e^{-s|E-\beta_j|^2}(\alpha_{\beta_j}X)e^{-s|E-\beta_j|^2}.
\]
Fourier transformation of the disjoint link kernels therefore
gives
\begin{equation}\label{eq:diag-mixed-trace}
 Z_T\left[\varphi_\beta\prod_aG_a\right]
 =(2\pi)^{|T|}\operatorname{Tr}_{H^l}
  \prod_{j=n_1-1,\ldots,0}
  \left(K_{\beta_j,S_j}^{1/2}B_jK_{\beta_j,S_j}^{1/2}
                         M_0u_jM_0^*\right),
\end{equation}
where
\[
 B_j=\prod_{a:j_a=j}(\alpha_{\beta_j}X_a)_\Xi,\qquad
 K_{\beta,S}=r_q^v
     e^{-|E-\beta|^2/2+|E-\beta|_S^2/4}P_G.
\]
Here $r_q^v=(2\pi/q)^{|\Sigma|}$ is the gauge-volume factor in
\eqref{eq:slab-electric-operators}; the primed kernels change
the electric Gaussian variance only on $S$. Both the kernels and the local operators
commute with $P_G$. In particular
$\|B_j\|\le\prod_{a:j_a=j}\|X_a\|$ by the norm identity above.

Set $T_j=K_{\beta_j,S_j}^{1/2}M_0$. These are
Hilbert--Schmidt by the same quotient Gaussian calculation as
\eqref{eq:slab-shifted-hs}. Move the first square root cyclically
to the end of the trace. The product then has alternating factors
$B_jT_ju_jT_{j-1}^*$. Schatten H\"older, with $2n_1$ equal
exponents, bounds its absolute trace by
\[
 \prod_j\|B_j\|\prod_j\|T_j\|_{2n_1}^2.
\]
The slice trace applied to $(T_j^*T_j)^{n_1}$ gives exactly
\begin{equation}\label{eq:diag-variance-trace}
 \operatorname{Tr}(T_j^*T_j)^{n_1}
 =(2\pi)^{-|T|}Z_T\left[
    \prod_{l=0}^{n_1-1}e^{i\langle\beta_j,F_l^{\mathrm{el}}\rangle}
    \prod_{l=0}^{n_1-1}\prod_{b\in S_j}
                     \sqrt2e^{-F_l^{\mathrm{el}}(b)^2/2}\right].
\end{equation}
The second product changes $g_1$ into $g_{1/2}$ on those links.
The positive bosonic law bounds the absolute value by
$2^{|S_j|n_1/2}(2\pi)^{-|T|}Z_T[1]$. Consequently
\[
 \prod_j\|T_j\|_{2n_1}^2
 \le 2^{\frac12\sum_j|S_j|}(2\pi)^{-|T|}Z_T[1].
\]
Together with \eqref{eq:diag-mixed-trace} this proves the two
operator bounds. For insertions in another time direction use
an orientation-preserving coordinate permutation. The scalar
bound follows directly from the positive law. Multiplication by
a character merely changes $\beta$, and enlargement restricts
the allowed placements. A signed permutation transports the
torus, the free columns and the placements; conjugation of a
character still has modulus one. These observations prove the
remaining assertions.
\end{proof}

\subsection{Amplification in two cell structures}\label{subsec:diag-amplification}

A diagonal fold exchanges the first and second time directions.
Its copies of a coordinate insertion therefore cannot all be put
in one coordinate trace. We will place these copies in different
coordinate cells and apply a second chessboard estimate. Each
resulting coordinate fold contains copies of only one insertion,
all with the same time direction, so
\eqref{eq:diag-coordinate-bound} applies.

The placement requires two scales depending on the support radius
$r$, but not its center or the tilt. Choose $h''\in h\Z$ and
$k\in\N$ such that
\begin{equation}\label{eq:diag-coarse-widths}
 h''\ge2r+6h,\qquad kh\ge16h'',\qquad
 kh\in h''\Z,\qquad4h'',2kh\in\mathcal D.
\end{equation}
These choices exist: take $4h''$ to be the least common multiple
of $4h$ and an element of $\mathcal D$ exceeding $4(2r+6h)$,
then take $2kh$ to be the least common multiple of $4h''$ and an
element exceeding $32h''$. Fix choices depending only on $h$, $r$
and the side sequence. For all sufficiently large $i$, the needed
divisibility conditions then hold simultaneously.

Here is the geometry of the placement. Let $B=y+[-r,r]^3$. Choose
coarse diagonal offsets $c'_\pm\in c_\pm+h\Z$ (with
$c_-=c_0$) and $c'_3\in c_3+h\Z$ so that
\begin{equation}\label{eq:diag-coarse-center}
 \ell_-(y)-c'_-\in[2h'',2h''+h),\quad
 \ell_+(y)-c'_+\in[4h'',4h''+h),\quad
 y_3-c'_3\in[h'',h''+h).
\end{equation}
The box lies strictly in the coarse cell $Q'_0$ of width $kh$.
Every element of its wall group has the form
\[
 g=\tau_{kh(z_-e_-+z_+e_++2z_3e_3)}
            \theta_-^{a_-}\theta_+^{a_+}\theta_3^{a_3},
 \qquad z\in\Z^3,\quad a\in\{0,1\}^3,
\]
where the reflections are at the lower walls of $Q'_0$. From
\eqref{eq:diag-coarse-center} one obtains
\begin{equation}\label{eq:diag-coordinate-cell-index}
 gy=y+h''\kappa(g)+s(a),\qquad |s(a)|_\infty<2h,
\end{equation}
with
\[
 \kappa(g)=(-2a_--4a_+,\,2a_--4a_+,\,-2a_3)
       +\frac{kh}{h''}(z_-+z_+,\,z_+-z_-,\,2z_3).
\]
The eight first vectors are distinct, and their pairwise
differences have size at most six. Since $kh/h''\ge16$, equality
of two indices first forces equality of their $a$'s and then of
their $z$'s. Thus $g\mapsto\kappa(g)$ is injective.

Choose half-integer coordinate offsets $o_\mu$ in
\[
 I_\mu=(y_\mu+r+2h-h'',\ y_\mu-r-2h).
\]
These intervals have length at least $2h$. Equation
\eqref{eq:diag-coordinate-cell-index} places every $gB$ strictly
inside coordinate cell $\kappa(g)$ of width $h''$. We can also
choose the offsets so that all these coordinate walls avoid
every reflected copy of $\supp O$. Under the original diagonal
wall group, the centers of the copies of a box containing $O$
have at most four residue classes modulo $h$ in either of the
first two coordinates, and at most two in the third. Its radius
$r_O$ therefore forbids at most $4(2r_O+1)=8r_O+4$ half-integer
classes. Since $h>8r_O+8$, an allowed class exists. Each $I_\mu$
contains a representative, and $h''\in h\Z$ keeps every wall
in the allowed class.

For large $i$, the adapted torus $D_i$ carries the original and
coarse diagonal structures; its period lattice lies in the
translation subgroup of the coarse wall group. The coordinate
structure of width $h''$ cuts $D_i$ and is carried by
$T_i=C(m_i/2,2m_i,m_i)$. Formula
\eqref{eq:diag-coordinate-cell-index} also shows that coarse
diagonal cells distinct modulo the periods have their boxes in
coordinate cells distinct modulo those periods. Indeed a period
translation changes $\kappa$ by that translation divided by
$h''$; injectivity then gives the converse. Thus the counts are
\begin{equation}\label{eq:diag-coarse-cell-counts}
 N'_i=\frac{m_i^3}{|Q|k^3},\qquad
 N''_i=\frac{m_i^3}{(h'')^3}.
\end{equation}
The first is the number of coarse diagonal cells in $D_i$; the
second is the number of coordinate cells in either torus. The
shortest periods are at least $m_i/2\ge kh\ge16h''$, so the
boxes, with the margins used for local lifts, embed. Folding
one such box in the coordinate structure gives a disjoint
placement with free columns, since every image stays inside
its open cell.

\begin{proposition}\label{thm:diag-amplification}
Put $\Delta_i=\Delta(m_i)$, with $\Delta$ defined in
\eqref{eq:slab-pressure-error}. It tends to zero as $i\to\infty$,
with the charge, mass and weights fixed.
If $G$ satisfies the coordinate bound $C$ in a box of radius $r$,
then, for $k$ as above and all sufficiently large $i$,
\begin{equation}\label{eq:diag-amplification}
 |\mu_i^t(G)|\le C\gamma_{t,i}(k),\qquad
 \gamma_{t,i}(k)=
      \exp(2t|O|_1k^3+|Q|k^3\Delta_i),\qquad t\ge0.
\end{equation}
The threshold and $k$ depend only on the side sequence, $h$ and
$r$. In particular they are uniform in the location of $G$.
\end{proposition}

\begin{proof}
Each coarse diagonal cell contains $k^3$ original cells. Put
$P=\exp(t\sum_{Q_\kappa\subset Q'_0}O^{[\kappa]})$. The tilt
$e^{tV_i}$ is exactly the product of the reflected copies of $P$
over the coarse cells: their walls are original walls, so each
reflection carries the original array to the same array. Apply
the adapted chessboard estimate to the integrands $PG$ in one
coarse cell and $P$ in the others. The fold of $P$ is $e^{tV_i}$.
After dividing its powers out, this gives
\begin{equation}\label{eq:diag-amplified-power}
 |\mu_i^t(G)|^{N'_i}
 \le\frac{Z_{D_i}[e^{tV_i}\prod_{a=1}^{N'_i}G^{g_a}]}
              {Z_{D_i}[e^{tV_i}]}.
\end{equation}
The right side is nonnegative by the reflected fold form. This
step uses reflection positivity of even integrands and does not
require $G$ itself to be a positive scalar function.

Expand $e^{tV_i}$ in curvature characters. The coefficient sum
is at most $\exp(tN_i|O|_1)$: the character algebra is closed
under multiplication, and the exponential series is absolutely
summable in its coefficient norm. Every plaquette in a character
of this expansion lies strictly within a coordinate cell,
because the chosen walls avoid the support of every copy of $O$.
For each character term, collect its factors by coordinate cells.
Each cell then contains a character alone or a character times
one $G^{g_a}$; precisely $N'_i$ cells have the latter form. Period
translations can be used to place the integrands in a chosen
coordinate brick without changing their value on $D_i$.

Use the chessboard estimate from the adapted to the coordinate
frame, \eqref{eq:slab-chessboard-coordinate}. A cell with a
character alone has fold bounded by $Z_{T_i}[1]$. A cell with
a copy of $G$ has a fold consisting of a character and its
coordinate images; the geometry above supplies the free columns
and disjoint embedding, so its fold is bounded by
$C^{N''_i}Z_{T_i}[1]$. Taking the chessboard roots and multiplying
the cell contributions gives, for every character term,
\[
 \left|Z_{D_i}\left[\varphi_\beta
                    \prod_{a=1}^{N'_i}G^{g_a}\right]\right|
                \le C^{N'_i}Z_{T_i}[1].
\]
The coefficient sum therefore bounds the numerator of
\eqref{eq:diag-amplified-power} by
$e^{tN_i|O|_1}C^{N'_i}Z_{T_i}[1]$. The denominator bound
\eqref{eq:diag-tilt-denominator} contributes a second
$e^{tN_i|O|_1}$. Finally, the pressure comparison gives
$Z_{T_i}[1]/Z_{D_i}[1]\le e^{m_i^3\Delta_i}$.
The root of order $N'_i$ gives
\eqref{eq:diag-amplification}.
\end{proof}

We also need exponential moments, for which a coordinate operator
bound is unnecessary. Repeat \eqref{eq:diag-amplified-power} with
$G=e^{\langle\beta,F\rangle}$ supported in $B$. It is an admissible
local integrand, since
$\langle\beta,F\rangle\le\|\beta\|^2+
\frac14\sum_{\supp\beta}F_\sigma^2$. Its diagonal copies have
disjoint supports. Remove the tilt from the numerator by its
pointwise bound, and apply the second chessboard estimate to
their coordinate cell integrands. The fold of one exponential
has $N''_i$ disjoint copies of $\beta$, and hence squared norm
$N''_i\|\beta\|^2$. Gaussian domination on $T_i$ bounds its
chessboard root by
$e^{\|\beta\|^2/2}Z_{T_i}[1]^{1/N''_i}$. There are $N'_i$
nonempty cells. The same denominator and pressure bounds
therefore yield
\begin{equation}\label{eq:diag-tilted-domination}
 \mu_i^t(e^{\langle\beta,F\rangle})
      \le e^{\|\beta\|^2/2}\gamma_{t,i}(k).
\end{equation}
For each fixed $t$ this proves tightness on $\Omega$. Indeed,
Markov's inequality applied to $\beta=\pm s1_\sigma$ gives a
uniform coordinate tail $2e^{-sR+s^2/2}\gamma_{t,i}(k(h,1))$.
Choose radii $R_\sigma$ with these tails summable over plaquettes;
the intersection of $\Omega$ with their product of compact
intervals is compact. The finitely many initial laws are tight
individually. Any weak limit $\nu_t$ consequently satisfies
\begin{equation}\label{eq:diag-limit-domination}
 \nu_t(e^{\langle\beta,F\rangle})
      \le e^{\|\beta\|^2/2+2t|O|_1k^3}.
\end{equation}
This uses lower semicontinuity of the exponential integral,
as explained after \eqref{eq:diag-tilted-positive-form}.

